% !TEX program = xelatex
% Overleaf: set compiler to XeLaTeX or LuaLaTeX.
\documentclass[12pt,a4paper]{ctexart}

\usepackage[a4paper,margin=2.5cm]{geometry}
\usepackage{booktabs}
\usepackage{longtable}
\usepackage{tabularx}
\usepackage{array}
\usepackage{float}
\usepackage{pdflscape}
\usepackage{xcolor}
\usepackage{hyperref}
\usepackage{enumitem}
\usepackage{listings}
\usepackage{caption}

\hypersetup{
  colorlinks=true,
  linkcolor=blue,
  citecolor=blue,
  urlcolor=blue
}

\setlist[itemize]{leftmargin=2em}
\setlist[enumerate]{leftmargin=2em}
\newcolumntype{Y}{>{\raggedright\arraybackslash}X}
\newcolumntype{R}{>{\raggedleft\arraybackslash}p{1.25cm}}
\renewcommand{\arraystretch}{1.18}

\lstdefinestyle{reportcode}{
  basicstyle=\ttfamily\small,
  breaklines=true,
  columns=fullflexible,
  frame=single,
  backgroundcolor=\color{gray!5}
}
\lstset{style=reportcode}

\title{LLM-Assisted Fault Localization 初步实验报告}
\author{Ziniu Jin}
\date{2026-05-30}

\begin{document}

\maketitle

\begin{abstract}
本报告总结一个基于检索增强和大语言模型重排序的文件级缺陷定位实验。系统首先使用 BM25 或 focused hybrid retrieval 生成候选文件，然后对不确定或存在领域错配的样本调用 LLM 进行选择性重排序。实验覆盖 Defects4J 六个 pilot 项目、AboutWork 公司项目和 Easy Finance 公司项目，并进一步测试 controlled agentic inspection 与 verifier pass。当前结果表明，selective LLM rerank 在高召回候选池上能显著提升 Top-k fault localization；agentic inspection 技术上可运行，但当前没有稳定超过 one-shot rerank；verifier pass 在 Easy Finance 上属于负向 ablation。
\end{abstract}

\tableofcontents
\newpage

\section{实验目标}

本实验研究一个文件级缺陷定位方法：给定 bug report、失败测试、stack trace 和源码候选文件，系统输出最可能包含缺陷的文件排序。

当前实验的核心问题不是自动修复 bug，而是回答：

\begin{lstlisting}
Where is the bug?
\end{lstlisting}

实验方法采用两阶段设计：

\begin{lstlisting}
BM25 / hybrid retrieval
-> candidate source files
-> selective LLM rerank
-> optional controlled agentic inspection
-> optional verifier
-> file-level ranking
\end{lstlisting}

ground truth 来自 verified fixing commit touched files，但 ground truth、patch、fix commit 内容不进入 prompt，只用于最终评价。

\section{研究问题}

本阶段实验主要对应表 \ref{tab:rq} 中的问题。

\begin{table}[H]
\centering
\caption{Research questions}
\label{tab:rq}
\begin{tabularx}{\textwidth}{@{}lY@{}}
\toprule
ID & Research Question \\
\midrule
RQ1 & 与 BM25 / hybrid retrieval 相比，LLM rerank 是否能提升文件级 fault localization？ \\
RQ2 & 候选池召回和 snippet evidence 质量如何影响最终结果？ \\
RQ3 & selective rerank 是否能在保持效果的同时降低 token 成本？ \\
RQ4 & 方法能否从 Defects4J benchmark 迁移到真实公司项目？ \\
RQ5 & controlled agentic inspection 和 verifier 是否比 one-shot rerank 更有效？ \\
\bottomrule
\end{tabularx}
\end{table}

\section{实验对象}

\subsection{Defects4J Benchmark}

Defects4J 当前完成 6 个 pilot 项目，每个项目 20 个 bug，共 120 个 bug：

\begin{lstlisting}
Lang-20
Math-20
Chart-20
Time-20
Closure-20
Mockito-20
\end{lstlisting}

这些项目主要用于验证方法在公开 Java benchmark 上的可行性。

\subsection{AboutWork 公司项目}

AboutWork 是真实公司项目 bug log 数据集：

\begin{lstlisting}
records: 39 committed-history bug logs
backend: 17
frontend: 22
source: /Users/jin/capi_project/aboutwork/COMPANY_BUG_LOG.md
\end{lstlisting}

它覆盖 Python/Django、TypeScript/React、chatbot/admin-agent 等场景，用于验证方法在真实产品项目上的表现。

\subsection{Easy Finance 公司项目}

Easy Finance 来自 backend/frontend git history：

\begin{lstlisting}
clean committed records: 63
backend: 29
frontend: 34
strict sensitivity set: 62
source: /Users/jin/capi_project/easy_finance/
\end{lstlisting}

Easy Finance 的 bug report 多数由 commit metadata 推断，因此比 AboutWork 手写 bug log 噪声更高。实验中额外构建了 strict62 sensitivity analysis，排除一条 bug text 与实际 touched files 明显不一致的样本。

\section{方法设计}

\subsection{Baseline Retrieval}

baseline 使用 BM25 或 focused hybrid retrieval 生成候选文件列表。query 主要由以下信息构成：

\begin{lstlisting}
bug_report.id
bug_report.text
test_failure
triggering_tests
stack_trace runtime context
\end{lstlisting}

BM25 的作用是提供候选池，而不是直接作为最终答案。后续 LLM rerank 只能在候选池内排序。

\subsection{LLM Rerank}

LLM rerank 接收：

\begin{itemize}
  \item bug report
  \item failing test / stack trace
  \item BM25 或 hybrid top-k candidate files
  \item file path、class/method metadata
  \item relevant source snippets
  \item optional retrieval evidence
\end{itemize}

输出为 JSON ranking。系统会过滤 invalid files、去重，并在模型返回不足时用原始 BM25 顺序 fallback。

\subsection{Selective Rerank}

为了控制成本，实验不对所有 bug 调用 LLM。selector 会根据低置信度、domain mismatch、management command noise、frontend/backend 业务词错配等规则选出 hard cases。

核心思想：

\begin{lstlisting}
Only call the LLM when retrieval is uncertain or likely wrong.
\end{lstlisting}

\subsection{Controlled Agentic Inspection}

为了回答 RQ5，实验加入 controlled agentic inspection。agent 不允许自由操作 repo，只能使用固定工具：

\begin{lstlisting}
search_files(query)
inspect_candidate(file)
read_file_window(file, start_line, end_line)
\end{lstlisting}

约束：

\begin{lstlisting}
candidate pool: BM25 top-50 only
max_steps: 2
no fix commit
no patch
no ground truth
final output: same prediction JSONL schema
\end{lstlisting}

\subsection{Verifier Pass}

verifier 是独立验证器，不允许再调用工具，只能看：

\begin{itemize}
  \item bug report
  \item agent top-10 ranking
  \item deterministic snippets
  \item prior agent observations
\end{itemize}

目标是测试 independent verification 是否能改进 agent ranking。

\section{评价指标}

主要指标见表 \ref{tab:metrics}。

\begin{table}[H]
\centering
\caption{Evaluation metrics}
\label{tab:metrics}
\begin{tabularx}{\textwidth}{@{}lY@{}}
\toprule
Metric & Meaning \\
\midrule
Top-1 Accuracy & 正确文件是否排第 1 \\
Top-3 Accuracy & 正确文件是否出现在前 3 \\
Top-5 Accuracy & 正确文件是否出现在前 5 \\
Top-10 Accuracy & 正确文件是否出现在前 10 \\
MRR & Mean Reciprocal Rank \\
Total Tokens & LLM 总 token 消耗 \\
Avg Tokens / Call & 平均每次调用 token 成本 \\
\bottomrule
\end{tabularx}
\end{table}

多文件 bug 使用 file-level hit：只要 ground truth files 中任意一个文件出现在 Top-k 内，即视为 Top-k 命中。

\section{Defects4J 结果}

当前 Defects4J best current row 汇总见表 \ref{tab:defects4j-results}。

\begin{landscape}
\begin{longtable}{@{}p{2.2cm}p{4.3cm}p{5.2cm}rrrr@{}}
\caption{Defects4J best current results}
\label{tab:defects4j-results}\\
\toprule
Project & Baseline / Retrieval & Best Current Method & LLM Calls & Top-1 & Top-3 & Top-5 & MRR \\
\midrule
\endfirsthead
\toprule
Project & Baseline / Retrieval & Best Current Method & LLM Calls & Top-1 & Top-3 & Top-5 & MRR \\
\midrule
\endhead
\bottomrule
\endfoot
Lang-20 & BM25 & BM25 + DeepSeek full rerank & 20 & 0.95 & 1.00 & 1.00 & 0.9750 \\
Math-20 & BM25 / hybrid & Hybrid/direct + compact evidence DeepSeek & 20 & 0.90 & 0.95 & 1.00 & 0.9350 \\
Chart-20 & Focused hybrid/direct & Focused hybrid/direct & 0 & 0.85 & 0.95 & 1.00 & 0.9040 \\
Time-20 & Focused hybrid/direct + test-prefix hints & Focused hybrid/direct + DeepSeek hard4 & 4 & 0.90 & 1.00 & 1.00 & 0.9500 \\
Closure-20 & Focused hybrid/direct & DeepSeek hard8 + pass-chain C13 + type-cycle C4 & 10 & 0.40 & 0.80 & 1.00 & 0.6225 \\
Mockito-20 & Focused hybrid/direct & Diagnostic hard7 + ByteBuddy M20 add-on & 8 & 0.55 & 0.90 & 1.00 & 0.7200 \\
\end{longtable}
\end{landscape}

Macro average:

\begin{lstlisting}
Top-1: 0.7583
Top-3: 0.9333
Top-5: 1.0000
MRR:   0.8511
\end{lstlisting}

初步解释：

\begin{itemize}
  \item 6 个 Defects4J pilot 项目均达到 Top-5 1.00。
  \item Lang、Math、Time 的结果较强，说明 LLM rerank 在高召回候选池上效果明显。
  \item Closure 和 Mockito 的 best row 包含 diagnostic add-on，因此当前更适合作为方法潜力证明，还不能直接声称完全自动化。
  \item 结果支持一个核心判断：LLM 适合作为 second-stage reranker，而不是替代 retrieval。
\end{itemize}

\section{AboutWork 结果}

AboutWork production-only 结果见表 \ref{tab:aboutwork-results}。

\begin{table}[H]
\centering
\caption{AboutWork results}
\label{tab:aboutwork-results}
\small
\begin{tabularx}{\textwidth}{@{}Yrrrrrr@{}}
\toprule
Method & LLM Calls & Top-1 & Top-3 & Top-5 & Top-10 & MRR \\
\midrule
BM25 production & 0 / 39 & 0.5897 & 0.8462 & 0.9487 & 0.9487 & 0.7171 \\
BM25 + selector\_v3 + DeepSeek & 11 / 39 & 0.7692 & 1.0000 & 1.0000 & 1.0000 & 0.8675 \\
\bottomrule
\end{tabularx}
\end{table}

提升：

\begin{lstlisting}
Top-1: 0.5897 -> 0.7692
Top-3: 0.8462 -> 1.0000
Top-5: 0.9487 -> 1.0000
MRR:   0.7171 -> 0.8675
\end{lstlisting}

意义：

\begin{itemize}
  \item AboutWork 证明 pipeline 可以迁移到真实公司项目。
  \item selective rerank 只调用 11/39 个样本，约 28.2\%，但 Top-3、Top-5、Top-10 均达到 1.00。
  \item 这说明真实项目中 LLM 不需要全量调用，关键是识别 BM25 可能犯错的样本。
\end{itemize}

限制：

\begin{itemize}
  \item selector\_v3 是在当前 39 条 bug log 上调出来的。
  \item 后续需要用新增 bug logs 验证泛化能力。
\end{itemize}

\section{Easy Finance 结果}

\subsection{Clean63 主结果}

Easy Finance clean63 production-only 结果见表 \ref{tab:easy-clean63}。

\begin{table}[H]
\centering
\caption{Easy Finance clean63 results}
\label{tab:easy-clean63}
\small
\begin{tabularx}{\textwidth}{@{}Yrrrrrr@{}}
\toprule
Method & Selected LLM Calls & Top-1 & Top-3 & Top-5 & Top-10 & MRR \\
\midrule
BM25 production & 0 / 63 & 0.3968 & 0.6825 & 0.8571 & 0.8730 & 0.5727 \\
BM25 + hard9 oracle & 9 / 63 & 0.4444 & 0.7778 & 0.9524 & 1.0000 & 0.6344 \\
BM25 + selector\_v1 & 10 / 63 & 0.4603 & 0.7778 & 0.9524 & 1.0000 & 0.6450 \\
BM25 + selector\_v1 UI evidence v2 & 10 / 63 & 0.4921 & 0.8095 & 0.9841 & 1.0000 & 0.6729 \\
\bottomrule
\end{tabularx}
\end{table}

当前 clean63 最好结果：

\begin{lstlisting}
BM25 production + selector_v1 UI evidence v2 + DeepSeek rerank
selected: 10 / 63
Top-1:  0.4921
Top-3:  0.8095
Top-5:  0.9841
Top-10: 1.0000
MRR:    0.6729
tokens: 250902 total
\end{lstlisting}

主要改进来自：

\begin{itemize}
  \item selector\_v1 捕捉 backend area mismatch、management command noise、frontend domain mismatch。
  \item UI evidence v2 强化 \texttt{useQuery}、button loading、currency display、unreviewed count 等代码证据。
  \item Top-10 从 0.8730 提升到 1.0000，说明候选召回和 selective rerank 对真实项目有效。
\end{itemize}

\subsection{Strict62 Sensitivity Analysis}

strict62 排除 \texttt{easyfinance-frontend-20250930-001}。该样本的 bug text 提到 confirm expense/invoice buttons，但实际 fix touches 涉及 Continue/Create Invoice/View PDF loading 状态，bug text 与 touched-file ground truth 存在明显不一致。

Strict62 结果见表 \ref{tab:easy-strict62}。

\begin{table}[H]
\centering
\caption{Easy Finance strict62 sensitivity analysis}
\label{tab:easy-strict62}
\small
\begin{tabularx}{\textwidth}{@{}Yrrrrrr@{}}
\toprule
Method & Selected LLM Calls & Top-1 & Top-3 & Top-5 & Top-10 & MRR \\
\midrule
BM25 production & 0 / 62 & 0.4032 & 0.6935 & 0.8710 & 0.8871 & 0.5810 \\
BM25 + selector\_v1 UI evidence v2 & 9 / 62 & 0.5000 & 0.8226 & 1.0000 & 1.0000 & 0.6817 \\
BM25 + selector\_v1 controlled agentic s2 & 9 / 62 & 0.5000 & 0.8065 & 1.0000 & 1.0000 & 0.6831 \\
BM25 + selector\_v1 agentic s2 + verifier & 9 / 62 & 0.5000 & 0.8065 & 1.0000 & 1.0000 & 0.6804 \\
\bottomrule
\end{tabularx}
\end{table}

解释：

\begin{itemize}
  \item strict62 上 one-shot UI evidence v2 已达到 Top-5 和 Top-10 1.00。
  \item controlled agentic s2 的 MRR 略高于 one-shot，但 Top-3 略低，且 token 成本更高。
  \item verifier pass 没有提升结果，反而略降 MRR。
\end{itemize}

\section{Agentic Inspection 与 Verifier 分析}

\subsection{Agentic Inspection}

Easy Finance strict62 selector\_v1 的 9 个样本上，agentic s2 单独结果：

\begin{lstlisting}
Top-1:  0.6667
Top-3:  0.8889
Top-5:  1.0000
Top-10: 1.0000
MRR:    0.8056
tokens: 250481
\end{lstlisting}

合并回 62 条后：

\begin{lstlisting}
Top-1:  0.5000
Top-3:  0.8065
Top-5:  1.0000
Top-10: 1.0000
MRR:    0.6831
\end{lstlisting}

初步结论：

\begin{itemize}
  \item controlled agent protocol 技术上可运行。
  \item 工具 trace 可记录，输出 JSON 可评价，能接入现有 merge/evaluation pipeline。
  \item 但当前 Easy Finance strict62 上，agentic inspection 没有明显超过 one-shot UI evidence rerank。
\end{itemize}

\subsection{Verifier Pass}

verifier 额外结果：

\begin{lstlisting}
selected 9-case verifier:
Top-1:  0.6667
Top-3:  0.8889
Top-5:  1.0000
MRR:    0.7870
tokens: 80569

merged strict62:
Top-1:  0.5000
Top-3:  0.8065
Top-5:  1.0000
MRR:    0.6804
\end{lstlisting}

与 agent-only 对比：

\begin{table}[H]
\centering
\caption{Agent-only versus agent with verifier}
\label{tab:agent-verifier}
\begin{tabularx}{\textwidth}{@{}Yrr@{}}
\toprule
Method & Full Strict62 MRR & Total Tokens \\
\midrule
Agentic s2 only & 0.6831 & 250481 \\
Agentic s2 + verifier & 0.6804 & 331050 \\
\bottomrule
\end{tabularx}
\end{table}

初步结论：

\begin{itemize}
  \item 当前 verifier-loop 是负向 ablation。
  \item independent verifier 没有提升 Top-k，反而因为 noisy top-10 evidence 把一个正确文件从 rank 2 降到 rank 4。
  \item 这说明更多 LLM 步骤不一定更好；证据选择质量比 agent step 数量更关键。
\end{itemize}

\section{成本分析}

当前主要 LLM usage 汇总见表 \ref{tab:cost}.

\begin{landscape}
\begin{longtable}{@{}p{6.0cm}rrrrr@{}}
\caption{LLM usage summary}
\label{tab:cost}\\
\toprule
Experiment & LLM Calls & Total Tokens & Avg Tokens / Call & Total Seconds & Avg Seconds / Call \\
\midrule
\endfirsthead
\toprule
Experiment & LLM Calls & Total Tokens & Avg Tokens / Call & Total Seconds & Avg Seconds / Call \\
\midrule
\endhead
\bottomrule
\endfoot
Math-20 full compact evidence & 20 & 536625 & 26831.3 & -- & -- \\
Math-20 selective rerank & 6 & about 188355 & about 31392.5 & -- & -- \\
Time-20 hard4 & 4 & about 117138 & 29284.5 & about 73.2 & 18.312 \\
Closure hard8 + C13 + C4 & 10 & 325602 & 32560.2 & about 193.5 & 19.351 \\
Mockito hard7 + M20 successful add-on & 8 & 177588 & 22198.5 & 167.729 & 20.966 \\
AboutWork selector\_v3 & 11 & 303537 & 27594.3 & -- & -- \\
Easy Finance clean63 selector\_v1 UI evidence v2 & 10 & 250902 & 25090.2 & 215.589 & 21.559 \\
Easy Finance strict62 selector\_v1 UI evidence v2 & 9 & 226860 & 25206.7 & 188.506 & 20.945 \\
Easy Finance strict62 agentic s2 & 9 & 250481 & 27831.2 & 154.399 & 17.155 \\
Easy Finance strict62 verifier pass & 9 & 80569 & 8952.1 & 121.980 & 13.553 \\
\end{longtable}
\end{landscape}

成本结论：

\begin{itemize}
  \item selective rerank 是必要的。全量 rerank 在中大型项目上 token 成本较高。
  \item agentic s2 比 one-shot strict62 多约 10.4\% tokens，但没有带来稳定收益。
  \item verifier 单次较便宜，但叠加在 agent 后总成本达到 331050 tokens，效果反而下降。
\end{itemize}

\section{初步结论}

\subsection{支持的结论}

当前结果支持以下判断：

\begin{enumerate}
  \item LLM-assisted fault localization 作为 second-stage reranker 是有效的。
  \item 高召回、低噪声候选池是效果上限。
  \item Prompt snippet / evidence 质量对 rerank 影响很大。
  \item Selective rerank 是成本控制的关键。
  \item 真实公司项目中，domain mismatch 和 retrieval noise 是主要错误来源之一。
  \item Agentic inspection 可以运行，但当前还没有证明比精心设计的 one-shot rerank 更好。
  \item Verifier-loop 当前是负向结果，说明 independent verification 需要更干净的 evidence selection。
\end{enumerate}

\subsection{对各 RQ 的初步回答}

\begin{table}[H]
\centering
\caption{Preliminary answers to research questions}
\label{tab:rq-answers}
\begin{tabularx}{\textwidth}{@{}lY@{}}
\toprule
RQ & 初步回答 \\
\midrule
RQ1 & LLM rerank 在 Defects4J 多个项目上显著提升 Top-k，尤其当候选池召回足够时。 \\
RQ2 & 候选池召回和 snippet evidence 是主要瓶颈；真实文件不在候选池中时 LLM 无法补救。 \\
RQ3 & selective rerank 有效；AboutWork 只调用 11/39，Easy Finance 只调用 10/63 或 9/62，也能显著提升 Top-k。 \\
RQ4 & 方法可以迁移到 AboutWork 和 Easy Finance，但公司数据质量、commit-derived reports 和 selector overfitting 需要单独讨论。 \\
RQ5 & controlled agent 可运行，但当前不优于 one-shot；verifier pass 没有提升，属于负向 ablation。 \\
\bottomrule
\end{tabularx}
\end{table}

\section{威胁与限制}

\subsection{Internal Validity}

\begin{itemize}
  \item ground truth 使用 fix commit touched files，可能包含伴随修改，而不一定都是 root cause。
  \item 多文件 bug 使用任一文件命中，可能高估实际定位能力。
  \item 一些 selector 和 evidence rule 是在当前样本上诊断后形成的，存在 fitted rule 风险。
  \item Closure 和 Mockito 的 best row 包含 diagnostic add-on，不能直接当作完全自动 selector 结果。
\end{itemize}

\subsection{External Validity}

\begin{itemize}
  \item Defects4J 主要是 Java benchmark，不能完全代表真实多语言系统。
  \item AboutWork 和 Easy Finance 都是公司项目 case study，样本量仍有限。
  \item Easy Finance 的 bug reports 来自 git commit metadata，文本质量弱于人工 bug log。
\end{itemize}

\subsection{Construct Validity}

\begin{itemize}
  \item 文件级定位比 method-level 或 line-level 更粗。
  \item Top-k accuracy 不能完全代表真实 debugging effort。
  \item token 成本、prompt 长度和模型版本变化会影响实际可用性。
\end{itemize}

\subsection{Reproducibility Risk}

\begin{itemize}
  \item DeepSeek 模型版本可能变化。
  \item Defects4J checkout、compile、test 对本地 JDK、缓存和环境有依赖。
  \item 公司项目 worktree 和历史 commit 数据需要固定记录。
\end{itemize}

\section{下一步工作}

短期优先级：

\begin{enumerate}
  \item 把 Defects4J、AboutWork、Easy Finance 结果统一成最终论文表格。
  \item 把 Closure / Mockito diagnostic add-on 改写为非 oracle selector 规则。
  \item 为 Easy Finance 和 AboutWork 增加新 bug logs，验证 selector 泛化能力。
  \item 对 agentic inspection 做更精准的 case selection，只在 one-shot snippet 明显不足的样本上测试。
  \item 暂时不把 verifier-loop 作为主方法，而是作为 negative ablation 写入 RQ5。
\end{enumerate}

论文写作建议：

\begin{itemize}
  \item 主结果应强调 BM25/hybrid + selective LLM rerank。
  \item Agent/verifier 不要写成主要贡献，而应写成 exploratory extension。
  \item 公司项目结果要单独标注数据来源差异：AboutWork 是手写 bug log，Easy Finance 是 git-history-derived reports。
  \item Threats to Validity 中必须说明 touched-file ground truth 的局限。
\end{itemize}

\section{主要结果文件}

主整理文档：

\begin{lstlisting}
project/fl-localizer/docs/results_integration_2026-05-27.md
project/fl-localizer/docs/stage_report_2026-05-27.md
project/fl-localizer/docs/current_results_report.md
project/fl-localizer/docs/worklog.md
\end{lstlisting}

Easy Finance agent/verifier 相关输出：

\begin{lstlisting}
project/fl-localizer/outputs/easy_finance_committed_strict62_agentic_deepseek_selector_v1_s2.jsonl
project/fl-localizer/outputs/easy_finance_committed_strict62_agentic_deepseek_selector_v1_s2_trace.jsonl
project/fl-localizer/outputs/easy_finance_committed_strict62_bm25_prod_plus_agentic_deepseek_selector_v1_s2_eval.json
project/fl-localizer/outputs/easy_finance_committed_strict62_agentic_plus_verifier_deepseek_selector_v1_s2.jsonl
project/fl-localizer/outputs/easy_finance_committed_strict62_bm25_prod_plus_agentic_verifier_deepseek_selector_v1_s2_eval.json
project/fl-localizer/outputs/easy_finance_committed_strict62_agentic_plus_verifier_deepseek_selector_v1_s2_usage.json
\end{lstlisting}

\section{一句话总结}

当前实验已经证明：在高召回候选池基础上，selective LLM rerank 能显著提升文件级 fault localization，并且能迁移到真实公司项目；但 agentic inspection 和 verifier-loop 目前还只是探索性扩展，尚未超过精心设计的 one-shot selective rerank。

\end{document}
